\begin{textsample}{POS Dim 3 – ro – Score 22.00 – ro/ro\_inf\_34.txt}  \label{ex:f3_pos_140}
Pensando nas \textbf{crianças} bem \textbf{pequenas} e \textbf{pequenas} - Propiciar a identificação e criação de diferentes \textbf{sons}, rimas e \textbf{gestos} em \textbf{brincadeiras} de \textbf{roda} e em outras interações sociais, ampliando sua linguagem oral. - Organizar momento para relatarem fatos acontecidos, histórias de \textbf{livros} que \textbf{ouviram}, ou assistiram na televisão filmes ou peças teatrais. - Oferecer situações de conversa com \textbf{adultos} e \textbf{crianças} sobre diferentes assuntos em diversos momentos da \textbf{rotina}. - Possibilitar a participação em \textbf{roda} de conversa onde expressam oralmente ideias, fatos, ou recontam histórias \textbf{escutadas}. - Propiciar situações de comunicação de regras básicas de alguns jogos aos \textbf{colegas}. - Oferecer situações de apreciação e comentários de leitura de histórias feita pelo(a) professor(a ). - Possibilitar a criação de histórias oralmente, a partir de imagens e temas sugeridos, bem como a recitação de parlendas e outros textos da tradição oral, como quadrinhas, adivinhas etc. - Oportunizar momento de revisitação das histórias nos \textbf{livros} que lhe são lidas e seus \textbf{personagens}. - Mediar uso de procedimentos básicos de um leitor, tais como ler a partir da capa e virar as páginas sucessivamente etc. - Possibilitar que imitem comportamentos de escritor ao escrever não convencionalmente recados e identificação da escrita do nome próprio em listas e objetos. - Propiciar situações em que as \textbf{crianças} possam relacionar texto e imagem e antecipar sentidos na leitura de quadrinhos, tirinhas e revistas de heróis. - Propiciar situações para reconhecerem o uso social de textos como : convites para festas de aniversário, roteiro de atividades do \textbf{dia}, comunicados aos pais e listas variadas. - Oportunizar situações de exploração de diferentes ferramentas e suportes de escrita para, a seu modo, desenhar, traçar letras e outros sinais gráficos. - Garantir a expressão na linguagem oral, musical e corporal, na dança, no desenho, na linguagem escrita, na dramatização e em outras linguagens em vários momentos. - Garantir a participação em \textbf{rodas} de conversa onde discutem seus pontos de vista sobre um assunto. - Oportunizar situações de debates de um assunto polêmico do cotidiano da unidade, por exemplo, como organizar o uso dos \textbf{brinquedos} do \textbf{parque}, etc. - Oportunizar a organização oral das etapas de uma tarefa, os passos de uma receita \textbf{culinária}, do preparo de uma tinta, ou as regras para uma \textbf{brincadeira}, por exemplo. - Oferecer situações de reconto de histórias a partir das narrativas da \textbf{professora} e/ou com apoio dos \textbf{livros}, utilizando recursos expressivos próprios e preservando os elementos da linguagem que se escreve. - Propiciar que exponham suas impressões sobre textos de prosa ou poesia que lhes foram lidos, bem como relatarem aos \textbf{colegas} histórias lidas por alguém de sua família. - Organizar situações de escolha e gravação de poemas para enviar a outras \textbf{crianças} ou aos familiares. - Oportunizar a participação em sarau literário, narrando ou recitando seus textos favoritos. - Possibilitar a criação de história de aventuras, definindo o ambiente em que ela ocorre, as características e desafios de seus \textbf{personagens}. - Possibilitar situação de documentação de reconto de história, tendo o professor como escriba. - Possibilitar \textbf{registros} escritos do nome sempre que for necessário para reconhecerem a semelhança entre a letra inicial de seu nome e as iniciais dos nomes dos \textbf{colegas} que possuem a mesma letra. - Possibilitar \textbf{registros} escritos de cartas, recados ou \textbf{diários} para determinada pessoa, elaboram convites, comunicados e listas, panfletos com as regras de um jogo, ainda que de modo não convencional. - Possibilitar que organizem com os \textbf{colegas}, com apoio da \textbf{professora}, coletâneas escritas de histórias, contos clássicos, lendas, contos populares, parlendas, \textbf{brincadeiras} cantadas, receitas \textbf{culinárias} etc. - Oportunizar que levantem hipóteses sobre o que está escrito e sobre como se escreve e utilizam conhecimentos sobre o sistema de escrita para localizar um nome específico em uma lista de palavras ( ingredientes de uma receita \textbf{culinária}, peças do jogo etc. ) ou palavras em um texto que sabem de memória. - Oportunizar exploração com os \textbf{colegas} de materiais impressos variados, de diferentes gêneros ( \textbf{livros} de literatura \textbf{infantil}, em verso e em prosa ; \textbf{livros} de imagem ; \textbf{livros} não ficcionais ; revistas ; jornais ; panfletos ; \textbf{embalagens} entre outros ).

% matched lemmas: adulto, brincadeira, brinquedo, colega, criança, culinário, dia, diário, embalagem, escutar, gesto, infantil, livro, ouvir, parque, pequeno, personagem, professora, registro, roda, rotina, som
\end{textsample}
